\chapter*{Abstract}
\addcontentsline{toc}{chapter}{Abstract}

Automated spike sorting allows researchers to cluster neural activity by neuron using electrophysiology, the recording and analysis of electrical activity in neurons. The most meaningful quality metric requires ground-truth data, which is rarely available during an experiment. This thesis examines whether two of the most meaningful quality metrics, \fpos{} (false positive rate of spikes) and \fmiss{} (false negative rate of spikes), can be predicted from features derived from the unit (putative neuron) in real paired juxtacellular and intracellular ground-truth recordings. We also explore the possibility of transfer learning from abundant hybrid simulation data, but the domain shift between hybrid simulation data and real paired ground-truth data is significant. On top of that, we investigate adding recording-context features, meaning summary features of the recording, to see whether they help predict these two targets.

To do so, this work creates a leakage-safe benchmark that evaluates the performance of different models on both \fmiss{} and \fpos{} across multiple evaluation protocols and recording families, from which the dataset is assembled.

The results indicate that the hybrid simulation data and paired ground-truth data are almost perfectly separable, emphasising the importance of the domain shift. The benchmark also revealed a structural asymmetry between the two targets, \fmiss{} and \fpos{}. Additionally, some models using transfer learning outperformed models that did not under specific settings. Furthermore, the context feature family was found to be particularly informative for \fpos{} prediction, while contributing much less consistently to \fmiss{}.

Overall, this thesis reveals that spike-sorting quality prediction is relevant for false positives but remains much more challenging for missed spikes.

\vspace{0.8em}
\noindent\textbf{Keywords:} spike sorting, electrophysiology, transfer learning, domain shift, quality control, gradient-boosted trees, SpikeForest
